\subsection{Patterns Worth Noticing}

When a matrix appears, inspect its structure before calculating. Useful questions
include:
\begin{itemize}
    \item Is the matrix diagonal or symmetric?
    \item Does it contain many zeros?
    \item Are two rows or columns equal?
    \item Is one row or column a scalar multiple of another?
    \item Does the matrix split naturally into simpler blocks?
\end{itemize}

\begin{interpretationbox}[Seeing structure is part of solving]
A visible pattern can explain why an operation simplifies, why two matrices
commute, or why some information is redundant. Recognising such a pattern is not
a computational trick added after the mathematics; it is part of understanding
the mathematical object.
\end{interpretationbox}

\begin{summarybox}[Read before calculating]
The zero, identity, transpose, symmetric, and diagonal patterns provide a first
vocabulary for structural reasoning. Before applying a general algorithm, check
whether the matrix already reveals a simpler route.
\end{summarybox}
